\section{Eliminating two obstruction degrees by Frobenius}
\label{frob:section}

Throughout this section $p=1009$.  We eliminate the rows $k=2$ and
$k=p-1$ of Table~\ref{may:remaining-table} in actual cobar cohomology.
The Frobenius extension reduces this task to quotient-group cohomology
with coefficients in kernel cohomology.  A filtration of those coefficients
then reduces the calculation to trivial-coefficient quotient groups, which
we compute by May methods.  We first establish these three steps and identify
the coefficient modules that the two obstruction degrees require.

\subsection{The Frobenius extension and its action}

Let $G_r$, for $r=2,3$, be the affine group scheme over
$\mathbb F_p$ with coordinate algebra
\[
 P_r=\mathbb F_p[t_1,\ldots,t_r],\qquad
 \Delta t_i=t_i\otimes1+1\otimes t_i+
       \sum_{0<l<i}t_l\otimes t_{i-l}^{p^l}.
\]
The weight of $t_i$ is $u_i$.  Thus, on points over a commutative
$\mathbb F_p$-algebra,
\begin{equation}\label{frob:group-law}
 (g h)_i=g_i+h_i+\sum_{0<l<i}g_lh_{i-l}^{p^l}.
\end{equation}
We write $H^a(G_r;M)$ for rational group cohomology, equivalently the
cohomology of the normalized coordinate cobar complex with coefficients
in the $G_r$-module $M$.

The spectral sequence below is the Lyndon--Hochschild--Serre spectral
sequence for an affine group scheme and a normal subgroup; see
\citet[Part~I, Proposition~6.6(3) and Section~6.7]{Jantzen1985} for
the spectral sequence and the quotient action on kernel cohomology.
We verify the required acyclicity and compute that action for this extension.

\begin{lemma}[The Frobenius extension]\label{frob:extension}
There is a finite faithfully flat exact sequence of affine group schemes
\begin{equation}\label{frob:exact-sequence}
 1\longrightarrow K_r=(\alpha_p)^r\longrightarrow G_r
   \xrightarrow{F} Q_r\longrightarrow1,
 \qquad O(Q_r)=\mathbb F_p[t_1^p,\ldots,t_r^p].
\end{equation}
Its derived invariants give a first-quadrant spectral sequence
\begin{equation}\label{frob:lhss}
 {}^{\mathrm F}E_2^{a,v}
   =H^a\bigl(Q_r;H^v(K_r;\mathbb F_p)\bigr)
   \Longrightarrow H^{a+v}(G_r;\mathbb F_p).
\end{equation}
It preserves weight and converges in every weight and total
cohomological degree.
\end{lemma}

\begin{proof}
The inclusion of the Frobenius subring makes $P_r$ free of rank $p^r$,
with basis $t_1^{e_1}\cdots t_r^{e_r}$, $0\leq e_i<p$.
It is therefore faithfully flat.  Its kernel has coordinate algebra
\[
 O(K_r)=\mathbb F_p[e_1,\ldots,e_r]/(e_1^p,\ldots,e_r^p).
\]
All the mixed terms in its coproduct vanish, so each $e_i$ is primitive.
This proves the description of the kernel and the exactness of
\eqref{frob:exact-sequence}.

Here is the acyclicity calculation needed to derive invariants.
Right multiplication by a kernel point is, by
\eqref{frob:group-law}, simply $t_i\mapsto t_i+e_i$: every mixed
term contains a positive $p$-power of a kernel coordinate.  Hence
multiplication gives an isomorphism of $K_r$-comodules
\begin{equation}\label{frob:regular-decomposition}
 O(Q_r)\otimes O(K_r)\xrightarrow{\ \cong\ }O(G_r),
 \qquad
 f\otimes e_1^{j_1}\cdots e_r^{j_r}
 \longmapsto f t_1^{j_1}\cdots t_r^{j_r},
 \quad 0\leq j_i<p.
\end{equation}
The first tensor factor has trivial $K_r$-coaction.  The second is the
regular comodule.  In particular,
$O(G_r)^{K_r}=O(Q_r)$, with its regular quotient coaction.

A cofree comodule is injective: its cofree functor is right adjoint to
the exact forgetful functor to vector spaces.  Equation
\eqref{frob:regular-decomposition} shows that every cofree
$G_r$-comodule is $K_r$-acyclic and that its $K_r$-invariants are
cofree over $Q_r$.  Every injective is a summand of a cofree comodule,
so these assertions also hold for injectives.  Finally,
$M^{G_r}=(M^{K_r})^{Q_r}$.  Resolving these two invariants functors
successively, using the indicated cofree resolutions, gives
\eqref{frob:lhss}.  This argument derives $K_r$-invariants; it does
not assert their exactness on arbitrary modules.  The spectral
sequence is first quadrant, so each fixed total cohomological degree
has a finite diagonal.  All coordinate weights are positive and each
weight component is finite dimensional.  These observations give the
asserted convergence, also weight by weight.
\end{proof}

\begin{lemma}[Conjugation on kernel cohomology]\label{frob:action}
For $r=3$ there is an isomorphism of bigraded algebras
\begin{equation}\label{frob:kernel-cohomology}
 H^*(K_3;\mathbb F_p)
   =\Lambda(x,a,c)\otimes\mathbb F_p[A,B,E].
\end{equation}
The exterior generators have cohomological degree one and weights
$1,u_2,u_3$, and the polynomial generators have cohomological degree
two and weights $p,pu_2,pu_3$, respectively.  Write
$\xi=t_1^p$, $\eta=t_2^p$ for quotient coordinates.  With the
right-conjugation convention, the quotient coaction is
\begin{align}
 x&\longmapsto x,&
 a&\longmapsto a+x\xi,&
 c&\longmapsto c+x\eta+a\xi^p,\label{frob:exterior-action}\\
 A&\longmapsto A,&
 B&\longmapsto B+A\xi^p,&
 E&\longmapsto E+A\eta^p+B\xi^{p^2}.
 \label{frob:polynomial-action}
\end{align}
Products in these formulas mean coefficient times quotient coordinate.
In particular, $xac$ and $A$ are invariant.  Exterior degree and
polynomial degree are separately preserved.  The corresponding
statements for $r=2$ omit $c$ and $E$.
\end{lemma}

\begin{proof}
The dual algebra of a primitive truncated coordinate
$\mathbb F_p[e]/(e^p)$ is a truncated polynomial algebra after rescaling
the dual basis by the invertible factorials $j!$, $j<p$.
Its resolution alternates multiplication by its generator and by the
$(p-1)$st power of that generator.  Applying the augmentation gives
one class in each cohomological degree.  Its algebra of extensions is
an exterior degree-one generator and a polynomial degree-two
generator.  The shift by two of the periodic resolution represents
the latter, and its iterates are nonzero in every even degree.
The former has square $[e\mid e]=\tfrac12d[e^2]$, which is zero
at this odd prime.
Tensoring these resolutions gives \eqref{frob:kernel-cohomology},
including the indicated weights.

For completeness, the naturality of the polynomial generators is
Frobenius semilinear.  A primitive coordinate $T$ has degree-two cocycle
\[
 \beta(T)=\sum_{j=1}^{p-1}\frac{1}{p}\binom pj
                      [T^j\mid T^{p-j}],
\]
where the displayed coefficients are integers reduced modulo $p$.
For a linear coordinate $T=\sum_i\lambda_i e_i$, the integral
polynomial
\[
 \frac{(\sum_i\lambda_i e_i)^p-\sum_i\lambda_i^p e_i^p}{p}
\]
has reduced coproduct
$\beta(\sum_i\lambda_i e_i)-\sum_i\lambda_i^p\beta(e_i)$.
Thus a matrix $(\lambda_{ij})$ on degree-one coordinates induces
$(\lambda_{ij}^p)$ on these degree-two cohomology classes.  This is a
polynomial identity in the $\lambda_{ij}$, so it also applies when
they are quotient coordinate functions.

If $e$ is a kernel point, direct use of \eqref{frob:group-law} gives
\[
 (g^{-1}eg)_1=e_1,\qquad
 (g^{-1}eg)_2=e_2+e_1g_1^p,\qquad
 (g^{-1}eg)_3=e_3+e_1g_2^p+e_2g_1^{p^2}.
\]
These coefficients belong to $O(Q_r)$, so the action descends to
the quotient.  The degree-one and degree-two naturality just proved
gives \eqref{frob:exterior-action} and
\eqref{frob:polynomial-action}.  The exterior matrix is triangular
with diagonal entries one, so its top exterior power $xac$ is fixed.
The formulas also prove separate preservation of the two degrees.
\end{proof}

\subsection{Coefficient filtration and quotient cohomology}

\begin{lemma}[The coefficient-weight filtration]\label{frob:coefficients}
Fix an exterior degree and a polynomial degree in
\eqref{frob:kernel-cohomology}, obtaining an invariant $Q_r$-submodule $M$.
At each total weight $w$, increasing coefficient weight filters the
normalized quotient cobar complex $C^*(Q_r;M)_w$ by a finite
filtration.  Its initial cohomology page is
\begin{equation}\label{frob:coefficient-page}
 {}^{\mathrm C}E_1^{a,\lambda}
      =M_{\lambda}\otimes H^a(Q_r;\mathbb F_p)_{w-\lambda},
\end{equation}
where $M_\lambda$ is viewed with trivial quotient action.
The coefficient differential with drop $d$ has the form
$ {}^{\mathrm C}d_d:(a,\lambda)\mapsto(a+1,\lambda-d)$.
If the nonidentity coefficient action has no weight drop smaller than
$d$, all earlier coefficient differentials vanish, and the differential
with drop $d$ is induced by that part of the coaction.

For the module $\operatorname{Sym}^n(A,B)$ we may write
\begin{equation}\label{frob:finite-D}
 A^n\langle1,D,\ldots,D^n\rangle,
 \qquad A^nD^j:=A^{n-j}B^j\quad(0\leq j\leq n).
\end{equation}
This is notation for a finite-dimensional module, with action
$D\mapsto D+\xi^p$; no inverse of $A$ is introduced.
\end{lemma}

\begin{proof}
In \eqref{frob:exterior-action} and
\eqref{frob:polynomial-action}, every nonidentity term has lower
coefficient weight, its difference being the positive weight of its
quotient coordinate.  The cobar differential therefore preserves
the increasing coefficient filtration.  In the associated graded
complex only the trivial coefficient action remains, giving
\eqref{frob:coefficient-page}.  At fixed total weight the coefficient
weights lie between zero and $w$, so the filtration is finite.
The stated description of its first possible differential follows by
taking the leading nonzero weight-lowering component of the cobar
differential.  Finally, expanding
\[
 (B+A\xi^p)^jA^{n-j}
   =\sum_{l=0}^j\binom jl A^{n-l}B^l\xi^{p(j-l)}
\]
proves the stated action on \eqref{frob:finite-D}.
\end{proof}

For either obstruction bidegree $(s,w)$, our aim is to make every group
$H^a(Q_r;H^{s-a}(K_r;\mathbb F_p))_w$ zero.  This makes the entire
total-degree-$s$ diagonal of \eqref{frob:lhss} zero, and hence proves
$H^{s,w}(G_r;\mathbb F_p)=0$.  We compute these groups by the coefficient
filtration \eqref{frob:coefficient-page}, using May only to bound or identify
its trivial-coefficient quotient groups.

The resulting vector-space upper bound has exactly the exterior-polynomial
letters used for $G_r$: kernel cohomology supplies the digit-zero letters,
and the quotient's initial May page supplies all higher digits.  This description
retains the kernel degree $v$.  Thus Table~\ref{may:remaining-table} bounds
every coefficient summand on the relevant Frobenius diagonal.
Lemma~\ref{may:comparison} reduces $k=2$ to $G_3$ and $k=p-1$ to $G_2$.
Grouping the candidates by their invariant kernel modules gives the four
computations in Table~\ref{frob:coefficient-obligations}.

For each module $M$ in that table, let $\lambda_0$ be its least coefficient
weight and put
\begin{equation}\label{frob:residual-weight}
 \mathcal R=\frac{w-\lambda_0}{p}.
\end{equation}
We call this the residual weight, and henceforth express quotient weights
after division by $p$.  Recall that $N_0=(p-1)^2$.
\begin{table}[htbp]
\centering
\begin{tabular}{ccccc}
\toprule
$k$ & $r$ & Kernel module $M$ & Quotient degree $a$ & $\mathcal R$\\
\midrule
$2$ & $3$ & $xac\operatorname{Sym}^{p-1}(A,B,E)$ & $2p-4$ & $p(p-1)$\\
$2$ & $3$ & $xac\operatorname{Sym}^{p-2}(A,B,E)$ & $2p-2$ & $p^2-p+1$\\
$p-1$ & $2$ & $\operatorname{Sym}^{N_0}(A,B)$ & $1$ & $3p+1$\\
$p-1$ & $2$ & $\operatorname{Sym}^{N_0-1}(A,B)$ & $3$ & $3p+2$\\
\bottomrule
\end{tabular}
\caption{The coefficient-cohomology groups that must vanish. Each row
means $H^a(Q_r;M)_w$, with $w=u_4+k$ in original weights.}
\label{frob:coefficient-obligations}
\end{table}

For the first two rows, the least weights are respectively
$\operatorname{wt}(xacA^{p-1})=3+p+2p^2$ and
$\operatorname{wt}(xacA^{p-2})=3+2p^2$; for the last two they are $np$,
where $n$ is the indicated polynomial degree.  These give the displayed
residual weights.  Replacing an $A$ by $E$ increases the coefficient weight,
after division by $p$, by $p^2+p$.  This exceeds both residual weights in
the $k=2$ rows, so no coefficient involving $E$ occurs in either fixed-weight
complex, in any quotient degree.  All four rows therefore use the finite
$D$-modules of \eqref{frob:finite-D}, multiplied by $xac$ in the first two
rows.  Suppressing the common invariant factor $A^n$ or $xacA^n$, a term
$D^j\zeta$ requires
\begin{equation}\label{frob:D-weight}
 \operatorname{wt}_Q(\zeta)=\mathcal R-jp.
\end{equation}
Their first possible coefficient differential lowers $j$ by one and raises
the quotient degree by one: its weight drop is $p^2$ in original weights,
or $p$ in these rescaled weights.

We now calculate the trivial-coefficient quotient groups needed for these
four computations and their adjacent degrees.
The coordinates $\xi,\eta$ have rescaled weights $1,p+1$ and
\begin{equation}\label{frob:quotient-coproduct}
 \overline\Delta\xi=0,\qquad
 \overline\Delta\eta=\xi\otimes\xi^p.
\end{equation}
In these rescaled weights, the letters
\[
 y=h_{1,0}^{Q},\quad b=h_{2,0}^{Q},\quad
 z=h_{1,1}^{Q},\quad C=b_{1,0}^{Q}
\]
have weights $1,p+1,p,p$, respectively; these are exactly the
letters $y,b,z,C$ used in Table~\ref{may:remaining-table}.

\begin{lemma}[The small quotient groups]\label{frob:quotient-groups}
For $Q_2$ or $Q_3$, at any positive rescaled weight
$w<p^2$ all available May letters are $y,b,z,C$.
Their monomials and bidegrees are
\begin{equation}\label{frob:small-monomials}
 C^j y^e z^f b^g,\quad
 j\geq0,\quad e,f,g\in\{0,1\},\qquad
 a=2j+e+f+g,\quad w=p(j+f+g)+e+g.
\end{equation}
The following lists are exhaustive, with a dash denoting the empty
list:
\begin{center}
\begin{tabular}{lll}
\toprule
Rescaled weight & Cohomological degree & May monomials\\
\midrule
$p(p-1)$ & $2p-4$ & ---\\
$p(p-1)$ & $2p-3$ & $zC^{p-2}$\\
$p(p-2)$ & $2p-4$ & $C^{p-2}$\\
$p(p-3)$ & $2p-5$ & ---\\
$p+2$ & $1,3$ & ---\\
$p+2$ & $2$ & $yb$\\
$2p+2$ & $2$ & ---\\
$2p+2$ & $3$ & $ybz$\\
$3p+2$ & $3$ & ---\\
$2$ & $1,3$ & ---\\
\bottomrule
\end{tabular}
\end{center}
The classes $C^{p-2}$ and $zC^{p-2}$ in the table are nonzero in actual
quotient cohomology.  There is also a nonzero class $\upsilon$ of degree two
and weight $p+2$, detected by $yb$, and $\upsilon z$ is nonzero
in degree three and weight $2p+2$.
\end{lemma}

\begin{proof}
The next exterior weights are $p^2$ and $p^2+p$; the third-coordinate
weight, if present, is $p^2+p+1$.  The next polynomial weight is
$p^2$.  All are outside the asserted range.  Lemma~\ref{may:page} therefore
gives exactly \eqref{frob:small-monomials}.  To verify the table,
first reduce its weight equation modulo $p$: its remainder is
$e+g\in\{0,1,2\}$.  For weights divisible by $p$ it forces
$e=g=0$, after which the degree equation fixes $j$ and $f$.
For the listed weights congruent to two it forces $e=g=1$;
substitution into both equations gives the entries shown.
The weight-two entries follow in the same way, since $j=f=g=0$
at that weight.  All weights in the table are within the asserted
range, because $3p+2<p(p-1)$ at $p=1009$.

The cocycles $z=[\xi^p]$ and $C=\beta(\xi)$ come from the
primitive coordinate $\xi$.  Their product $zC^{p-2}$ is an
actual cocycle with the indicated nonzero leading May monomial.
The preceding cohomological degree at its weight has no May
monomial in any filtration.  It consequently cannot become a
boundary on any May page.  An actual cocycle has no outgoing May
differential, proving its nonvanishing.  The factor $C^{p-2}$ is therefore
nonzero as well.

For the other two classes, use the cobar convention
$d[u\mid v]=d[u]\mid v-u\mid d[v]$ with
$d[\eta]=[\xi\mid\xi^p]$.  The corrected cocycle
\begin{equation}\label{frob:corrected-cocycle}
 U=[\xi\mid\eta]+\frac12[\xi^2\mid\xi^p]
\end{equation}
is closed: its two differentials are
$-[\xi\mid\xi\mid\xi^p]$ and
$[\xi\mid\xi\mid\xi^p]$, respectively.  In the ordinary
May filtration its first summand has degree four and the correction
has degree three.  Its leading detector is thus $yb$.  The table
excludes both an incoming and an outgoing May differential, so
$\upsilon=[U]\neq0$.  Similarly
\[
 U\mid\xi^p
   =[\xi\mid\eta\mid\xi^p]
       +\frac12[\xi^2\mid\xi^p\mid\xi^p]
\]
is a closed representative with leading detector $ybz$.
Its correction has smaller May filtration.  The absence of a
degree-two May monomial at weight $2p+2$ excludes every possible
incoming differential, and the explicit cocycle excludes outgoing
differentials.  Hence $\upsilon z\neq0$.
\end{proof}

\subsection{The row \texorpdfstring{$k=2$}{k=2}}

\begin{proposition}[The row $k=2$]\label{frob:k2}
At $s=4p-3$ and $w=u_4+2$, the cobar cohomology group in
Table~\ref{may:remaining-table} is zero.
\end{proposition}

\begin{proof}
We prove vanishing for $G_3$ by treating the first two rows of
Table~\ref{frob:coefficient-obligations}.  Their kernel degrees are
$2p+1$ and $2p-1$, so their coefficient differentials cannot mix.

For $M=xac\operatorname{Sym}^{p-1}(A,B,E)$, we have
$\mathcal R=p(p-1)$ and quotient degree $a=2p-4$.
After the weight exclusion of $E$, the finite module is
$xacA^{p-1}\langle1,D,\ldots,D^{p-1}\rangle$.
The sole candidate from Table~\ref{may:remaining-table} is
$D C^{p-2}$, suppressing the invariant factor $xacA^{p-1}$.
Its quotient weight is $p(p-2)$ by \eqref{frob:D-weight}.
The translation of $D$ gives, on the first possible coefficient page,
\begin{equation}\label{frob:k2-differential}
 D[C^{p-2}]\longmapsto \pm[zC^{p-2}].
\end{equation}
The target has quotient degree $2p-3$ and weight $p(p-1)$.
Both source and target are nonzero actual quotient-cohomology classes
by Lemma~\ref{frob:quotient-groups}; they are present on this page because
all earlier coefficient differentials vanish.  Thus
\eqref{frob:k2-differential} removes the candidate.

We record the adjacent-degree checks.  A differential on this page
entering its source would have $D$-degree two, quotient degree
$2p-5$, and quotient weight $p(p-3)$; this group is zero
by Lemma~\ref{frob:quotient-groups}.  The target, having $D$-degree
zero, has no outgoing coefficient differential.  Its first-page
incoming source can only have $D$-degree one, and
\eqref{frob:small-monomials} in degree $2p-4$, weight $p(p-2)$,
gives only $C^{p-2}$.  Thus no other source can cancel the nonzero
map in \eqref{frob:k2-differential}.  In the original total degree,
the same equation, or the exhaustive candidate table, leaves no
other coefficient-leading term in this vertical summand.  Hence
a change of representative cannot replace the removed class.

Now take the full second summand
$M=xac\operatorname{Sym}^{p-2}(A,B,E)$, with $\mathcal R=p^2-p+1$
and quotient degree $2p-2$.
Table~\ref{may:remaining-table} leaves only its $D$-degree-zero candidate,
whose coefficient $xacA^{p-2}$ is invariant.
Its quotient detector $yzC^{p-2}$ is already a
splitting boundary, since $db=yz$ and $dC=0$ on the May page.
This also has a direct cobar realization: $yz=d[\eta]$ and $C$
is an actual cocycle, so $yzC^{p-2}$ is exact.  It is the only
May monomial in this quotient bidegree, as follows by putting
$a=2p-2$, $w=p^2-p+1$ in \eqref{frob:small-monomials}.
It therefore contributes no actual quotient cohomology in
\eqref{frob:coefficient-page}.  These are all initial candidates
of total degree $4p-3$ and weight $u_4+2$.  The coefficient spectral
sequences consequently vanish there, so all corresponding terms
of \eqref{frob:lhss} are zero.  Its convergence proves the result.
\end{proof}

\subsection{The row \texorpdfstring{$k=p-1$}{k=p-1}}

\begin{proposition}[The row $k=p-1$]\label{frob:kpminus1}
Put $N_0=(p-1)^2$.  At $s=2N_0+1$ and
$w=u_4+p-1$, the cobar cohomology group in
Table~\ref{may:remaining-table} is zero.
\end{proposition}

\begin{proof}
We treat the last two rows of Table~\ref{frob:coefficient-obligations}
for $G_2$, using the finite module \eqref{frob:finite-D} in each case.
For $n=N_0$, the residual weight is $\mathcal R=3p+1$ and the quotient
degree is one.  We show directly that every one-cocycle vanishes.
The leading candidates could survive only if lower powers of $D$
supplied corrections.  The crossed-cocycle equation will rule these out:
it requires unequal mixed coefficients where every available correction
has equal ones.

Write $\xi(g)=r$, $\eta(g)=u$,
$\xi(h)=s$, $\eta(h)=v$, so that
\[
 \xi(gh)=r+s,\qquad\eta(gh)=u+v+rs^p.
\]
A normalized one-cocycle is a polynomial $f(D;g)$ satisfying
\begin{equation}\label{frob:crossed-equation}
 f(D;gh)=f(D;g)+f(D+r^p;h).
\end{equation}
Since $D$ has weight $p$ and $4p>3p+1$, it must have the form
\[
 f(D;g)=D^3f_3(g)+D^2f_2(g)+Df_1(g)+f_0(g).
\]
All these powers are in the finite module, since $N_0\geq3$.
For a polynomial function $\varphi$ write
$\delta\varphi(g,h)=\varphi(gh)-\varphi(g)-\varphi(h)$.
Comparing the four coefficients in
\eqref{frob:crossed-equation} gives
\begin{align}
 \delta f_3&=0,\nonumber\\
 \delta f_2&=3r^pf_3(h),\nonumber\\
 \delta f_1&=2r^pf_2(h)+3r^{2p}f_3(h),
       \label{frob:crossed-coefficients}\\
 \delta f_0&=r^pf_1(h)+r^{2p}f_2(h)+r^{3p}f_3(h).
       \nonumber
\end{align}
Homogeneity forces
\begin{align*}
 f_3&=a_3\xi,\\
 f_2&=b_2\eta+c_2\xi^{p+1},\\
 f_1&=d_1\eta\xi^p+e_1\xi^{2p+1},\\
 f_0&=r_0\xi^{3p+1}+s_0\eta\xi^{2p}
                       +t_0\eta^2\xi^{p-1}.
\end{align*}
Here the coefficients $a_3,b_2,c_2,d_1,e_1,r_0,s_0,t_0$ lie in $\mathbb F_p$.
The second equation of \eqref{frob:crossed-coefficients}, together
with
\[
 \delta\eta=rs^p,\qquad
 \delta(\xi^{p+1})=r^ps+rs^p,
\]
forces $c_2=3a_3$ and $b_2=-3a_3$.
The next equation has coefficient $2b_2$ on $r^pv$ and zero
on $us^p$.  On the other hand,
\begin{equation}\label{frob:asymmetric-pair}
 \delta(\eta\xi^p)
       =us^p+vr^p+r^{p+1}s^p+rs^{2p},
\end{equation}
and $\delta(\xi^{2p+1})$ contains neither $u$ nor $v$.
Its left side must therefore have equal coefficients on $r^pv$
and $us^p$.  Since $2b_2=-6a_3$ and $6\neq0$ in
$\mathbb F_{1009}$, we obtain $a_3=0$, and then
$b_2=c_2=0$.  No term in $f_0$ can alter this
coefficient-of-$D$ equation.

Now $f_1$ is primitive.  Equation~\eqref{frob:asymmetric-pair}
first forces $d_1=0$, and the coefficient of $r^{2p}s$
in $\delta(\xi^{2p+1})$ then forces $e_1=0$.
Finally $f_0$ is primitive.  In its reduced coproduct the monomial
$u^2s^{p-1}$ has coefficient $t_0$, so $t_0=0$.
The monomial $us^{2p}$ then has coefficient $s_0$, so
$s_0=0$.  The coefficient of $r^{3p}s$ in
$\delta(\xi^{3p+1})$ is one, giving $r_0=0$.
Every normalized cocycle is thus zero, proving the required
$H^1$ vanishing, including all lower-$D$ corrections.

For $n=N_0-1$, the residual weight is $\mathcal R=3p+2$ and the desired
quotient degree is three.  The first coefficient differential
comes from the term $2D\xi^p$ in
\[
 (D+\xi^p)^2-D^2=2D\xi^p+\xi^{2p}.
\]
Using the actual class $\upsilon$ of
Lemma~\ref{frob:quotient-groups}, it is
\begin{equation}\label{frob:h3-differential}
 D^2\upsilon\longmapsto\pm2D\upsilon z.
\end{equation}
The coefficient differential first prepends $z$. Its product
$z\upsilon$ agrees with $\upsilon z$: both are actual cocycles
with the same leading May monomial $ybz$, with coefficient $+1$
after two exterior interchanges, and the list in Equation~\eqref{frob:small-monomials}
leaves no other initial May monomial or lower-filtration cohomology
in this bidegree. Thus the product order in
Equation~\eqref{frob:h3-differential} is justified directly.
Both powers of $D$ lie in the finite module, since $n\geq3$.
Equation~\eqref{frob:D-weight} gives quotient degree and weight
$(2,p+2)$ for the source and $(3,2p+2)$ for the target.
Both therefore have the required residual weight $3p+2$.
All earlier coefficient differentials vanish, so source and target are
still their actual quotient-cohomology classes on this first possible
page.  They are both nonzero by
Lemma~\ref{frob:quotient-groups}.

A same-page incoming differential to the source would require
$D$-degree three and quotient $H^1$ of weight two,
which is zero.  The only same-page incoming source for the target
is $D^2H^2(Q_2)_{p+2}$, the one-dimensional span displayed in
\eqref{frob:h3-differential}.  A same-page outgoing differential
from the target is proportional to $\upsilon z^2=0$.
Indeed $[\xi^p\mid\xi^p]=\tfrac12d[\xi^{2p}]$, so this
last vanishing also follows directly in the cobar complex.
The term $\xi^{2p}$ in the expansion of $D^2$ lowers two
$D$-levels and cannot cancel its first differential.  Since
$2\neq0$ in $\mathbb F_p$, \eqref{frob:h3-differential}
removes the target.

There is no other degree-three coefficient-leading term.  In fact,
the possibilities for $D$-degrees $0,1,2,3$ require quotient
degree three at respective weights
$3p+2,2p+2,p+2,2$; by
Lemma~\ref{frob:quotient-groups}, only the second has a
candidate, namely $ybz$.  Larger $D$-degrees exceed the residual
weight.  Thus there is no replacement in a lower coefficient
filtration.  The initial candidate table excludes every other
kernel degree or exterior summand at this total bidegree.
Both possible terms of \eqref{frob:lhss} consequently vanish,
and the proposition follows.
\end{proof}

\begin{corollary}\label{frob:vanishing}
The obstruction groups in bidegrees
\[
 (s,w)=(2q+1,u_4+2),\qquad
 (s,w)=(q(p-1)+1,u_4+p-1)
\]
vanish in the additive cobar calculation of
Table~\ref{may:remaining-table}.  Among its rows with $k\geq2$,
only the row $k=p$ remains to be addressed.
\end{corollary}

\begin{proof}
Apply Propositions~\ref{frob:k2} and~\ref{frob:kpminus1},
and the exhaustive list in Table~\ref{may:remaining-table}.
\end{proof}
